\documentclass{labreport}

\usepackage{microtype}
\usepackage{enumitem}
\usepackage{pifont}
\usepackage{tabularx}
\usepackage{booktabs}

% Headheight adjustment for fancyhdr
\setlength{\headheight}{28pt}
\addtolength{\topmargin}{-13pt}

% Colors
\definecolor{codebg}{RGB}{248, 249, 251}
\definecolor{codeborder}{RGB}{218, 222, 229}
\definecolor{linkcolor}{RGB}{20, 60, 120}
\definecolor{zshgreen}{RGB}{0, 160, 80}
\definecolor{zshred}{RGB}{200, 30, 30}

% Hyperref setup with subtle link colors
\hypersetup{
    colorlinks=true,
    linkcolor=linkcolor,
    filecolor=linkcolor,
    urlcolor=linkcolor,
    citecolor=linkcolor
}

% Code and terminal listings
\lstdefinestyle{bash}{
    basicstyle=\ttfamily\footnotesize,
    backgroundcolor=\color{codebg},
    frame=single,
    framerule=0.5pt,
    rulecolor=\color{codeborder},
    breaklines=true,
    breakatwhitespace=false,
    showstringspaces=false,
    keepspaces=true,
    columns=flexible,
    numbers=none,
    xleftmargin=3pt,
    xrightmargin=3pt,
    framesep=4pt
}

\lstdefinestyle{c}{
    language=C,
    basicstyle=\ttfamily\footnotesize,
    backgroundcolor=\color{codebg},
    frame=single,
    framerule=0.5pt,
    rulecolor=\color{codeborder},
    keywordstyle=\color{blue!85!black}\bfseries,
    commentstyle=\color{gray!80!black}\itshape,
    stringstyle=\color{teal!80!black},
    numbers=left,
    numberstyle=\tiny\color{gray!80!black},
    numbersep=8pt,
    breaklines=true,
    breakatwhitespace=false,
    showstringspaces=false,
    keepspaces=true,
    columns=flexible,
    tabsize=2,
    xleftmargin=14pt,
    xrightmargin=3pt,
    framesep=4pt
}

\lstset{
    extendedchars=true,
    literate=
        {➜}{{\textcolor{zshgreen}{\ding{220}}}}{1}
        {├─}{|-\,}{2}
        {└─}{\`-\,}{2}
        {├}{|}{1}
        {─}{-}{1}
        {└}{\`}{1}
        {✗}{{\textcolor{zshred}{\ding{55}}}}{1}
}

% Section styling for clean document structure
\titleformat{\section}
  {\large\bfseries}{\thesection}{0.8em}{}[\vspace{1pt}\hrule height 0.5pt]
\titleformat{\subsection}
  {\normalsize\bfseries}{\thesubsection}{0.8em}{}
\titleformat{\subsubsection}
  {\small\bfseries}{\thesubsubsection}{0.8em}{}

\begin{document}

\labtitle
    {10}
    {Add Your Own System Call: trace}
    {st. gr. FAF-243}
    {Șoimu Ionuț}
    {asist. univ.\\[0.15cm] Reițman Patricia}

\tableofcontents
\newpage

% ==============================================================================
\section{Lab 10: Track B -- Add Your Own System Call: trace (xv6)}

\subsection{Before you start}

\subsubsection*{Purpose and Scope}
In Laboratory Work 9, we examined system calls from the user-space perspective on Linux using the external diagnostic utility \texttt{strace}. In Laboratory Work 10, the focus shifts directly into \textbf{kernel space}: we modify the kernel of \textbf{xv6-riscv} to implement our own native tracing mechanism from within.

Specifically, we design and add a new system call, \texttt{trace(int mask)}, that configures a calling process (and all of its future child processes) to intercept and print diagnostic messages whenever targeted system calls execute. By tracking syscall index \textbf{23} from a user application down through assembly stubs, the RISC-V \texttt{ecall} hardware exception mechanism, the kernel trap vector, the dispatch table, and the process control block, we gain a comprehensive understanding of the complete lifecycle of a system call.

\subsubsection*{Architectural Nuances of our xv6 Tree}
In older iterations of the MIT 6.1810 educational operating system course, syscall 22 is commonly assigned to \texttt{SYS\_trace}, and the kernel outputs diagnostic text using \texttt{printf()}. In our updated course tree:
\begin{itemize}[noitemsep]
    \item System call index \texttt{22} is already reserved for \texttt{SYS\_sync};
    \item We assign index \textbf{\texttt{23}} to \textbf{\texttt{SYS\_trace}};
    \item The kernel-space print routine is \textbf{\texttt{printk()}} (reserving \texttt{printf()} for user-space standard I/O).
\end{itemize}

% ==============================================================================
\subsection{Part 1. The number and the user side}

\subsubsection{Overview of User-Side Components}
Before the kernel can handle a new request, both the build system and user-space binaries must be made aware of the new system call signature and index number. This requires coordinating five specific files:

\begin{enumerate}[leftmargin=*]
    \item \textbf{1.1 System Call Number Definition (\texttt{kernel/syscall.h}):}\\
    We assign integer \texttt{23} to identify the new system call:
\begin{lstlisting}[style=c]
#define SYS_sync   22
#define SYS_trace  23
\end{lstlisting}

    \item \textbf{1.2 User Header Declaration (\texttt{user/user.h}):}\\
    We expose the C function prototype to user applications:
\begin{lstlisting}[style=c]
int sync(void);
int trace(int);
\end{lstlisting}

    \item \textbf{1.3 Assembly Stub Generation (\texttt{user/usys.pl}):}\\
    In xv6, user-space system call stubs are generated at build time by the Perl script \texttt{usys.pl}, which outputs the assembly file \texttt{user/usys.S}. We append the entry definition:
\begin{lstlisting}[style=bash]
entry("sync");
entry("trace");
\end{lstlisting}
    When compiled, this generates the machine instruction sequence:
\begin{lstlisting}[style=bash]
.global trace
trace:
 li a7, SYS_trace
 ecall
 ret
\end{lstlisting}
    The stub places syscall number \texttt{23} into register \texttt{a7}, places the argument \texttt{mask} into register \texttt{a0}, and executes \texttt{ecall} to trap into supervisor mode.

    \item \textbf{1.4 User Trace Utility (\texttt{user/trace.c}):}\\
    We implement the user-facing command-line utility \texttt{trace}, which accepts a bitmask integer, calls \texttt{trace(mask)} to configure its own process state, and executes the target program with \texttt{exec()}:
\begin{lstlisting}[style=c]
#include "kernel/types.h"
#include "kernel/stat.h"
#include "user/user.h"

int
main(int argc, char *argv[])
{
  if (argc < 3) {
    fprintf(2, "usage: trace mask command [args...]\n");
    exit(1);
  }
  if (trace(atoi(argv[1])) < 0) {
    fprintf(2, "trace: failed\n");
    exit(1);
  }
  exec(argv[2], argv + 2);
  fprintf(2, "trace: exec %s failed\n", argv[2]);
  exit(1);
}
\end{lstlisting}
    Because \texttt{exec()} replaces the process memory space without changing the process control block (\texttt{struct proc}), the \texttt{trace\_mask} set by \texttt{trace} persists into the newly loaded program image.

    \item \textbf{1.5 Build Integration (\texttt{Makefile}):}\\
    We append \texttt{\$U/\_trace\textbackslash} to the \texttt{UPROGS} target list:
\begin{lstlisting}[style=bash]
	$U/_wc\
	$U/_zombie\
	$U/_trace\
\end{lstlisting}
\end{enumerate}

\subsubsection{Check 1: Verification of User-to-Kernel Number Dispatch}
If we compile and run \texttt{trace} in QEMU before providing a kernel implementation, the user stub correctly executes \texttt{ecall} with \texttt{a7 = 23}. Because index 23 is not yet populated in the kernel dispatch table, the kernel trap handler reports an unknown system call:
\begin{lstlisting}[style=bash]
$ trace 32 grep hello README
3 trace: unknown sys call 23
trace: failed
\end{lstlisting}
This diagnostic message confirms that hardware register \texttt{a7} successfully conveys number 23 across the user/supervisor hardware boundary.

% ==============================================================================
\subsection{Part 2. The kernel side}

\subsubsection{Extending Kernel Data Structures and Dispatch Logic}
To store, manage, and execute the tracing logic inside the kernel, we modify four core kernel files:

\begin{enumerate}[leftmargin=*]
    \item \textbf{2.1 Process Control Block (\texttt{kernel/proc.h}):}\\
    Each active process must store its personal tracing bitmask. We extend \texttt{struct proc} by adding \texttt{trace\_mask}:
\begin{lstlisting}[style=c]
  char name[16];               // Process name (debugging)
  int trace_mask;              // which calls to print
};
\end{lstlisting}

    \item \textbf{2.2 Kernel Handler Implementation (\texttt{kernel/sysproc.c}):}\\
    We implement \texttt{sys\_trace()}, which reads argument 0 (passed in register \texttt{a0}) using \texttt{argint()} and saves it into the current process's \texttt{trace\_mask}:
\begin{lstlisting}[style=c]
uint64
sys_trace(void)
{
  int mask;
  argint(0, &mask);
  myproc()->trace_mask = mask;
  return 0;
}
\end{lstlisting}

    \item \textbf{2.3 Syscall Table Registration (\texttt{kernel/syscall.c}):}\\
    We declare the external prototype and insert the handler into the dispatch array:
\begin{lstlisting}[style=c]
extern uint64 sys_trace(void);
...
static uint64 (*syscalls[])(void) = {
...
  [SYS_sync]    = sys_sync,
  [SYS_trace]   = sys_trace,
};
\end{lstlisting}

    \item \textbf{2.4 Name Table and Interception Logic (\texttt{kernel/syscall.c}):}\\
    The kernel identifies calls internally by numeric constants. To display human-readable strings, we define an array of syscall names indexed by number placed immediately above \texttt{syscall()}:
\begin{lstlisting}[style=c]
static char *syscallnames[] = {
  [SYS_fork]    "fork",
  [SYS_exit]    "exit",
  [SYS_wait]    "wait",
  [SYS_pipe]    "pipe",
  [SYS_read]    "read",
  [SYS_kill]    "kill",
  [SYS_exec]    "exec",
  [SYS_fstat]   "fstat",
  [SYS_chdir]   "chdir",
  [SYS_dup]     "dup",
  [SYS_getpid]  "getpid",
  [SYS_sbrk]    "sbrk",
  [SYS_pause]   "pause",
  [SYS_uptime]  "uptime",
  [SYS_open]    "open",
  [SYS_write]   "write",
  [SYS_mknod]   "mknod",
  [SYS_unlink]  "unlink",
  [SYS_link]    "link",
  [SYS_mkdir]   "mkdir",
  [SYS_close]   "close",
  [SYS_sync]    "sync",
  [SYS_trace]   "trace",
};
\end{lstlisting}
    Inside \texttt{syscall()}, immediately following the execution of the call, we test whether bit \texttt{num} is enabled in \texttt{p->trace\_mask}. If enabled, we print the PID, syscall name, and return code (stored in \texttt{p->trapframe->a0}):
\begin{lstlisting}[style=c]
    p->trapframe->a0 = syscalls[num]();
    if ((p->trace_mask >> num) & 1)
      printk("%d: syscall %s -> %ld\n",
             p->pid, syscallnames[num], p->trapframe->a0);
\end{lstlisting}

    \item \textbf{2.5 Mask Inheritance Across Fork (\texttt{kernel/proc.c}):}\\
    In Unix, tracing utilities must track child processes spawned by the target application. Inside \texttt{fork()}, we ensure the child process \texttt{np} inherits its parent's mask:
\begin{lstlisting}[style=c]
  safestrcpy(np->name, p->name, sizeof(p->name));
  np->trace_mask = p->trace_mask;
\end{lstlisting}
\end{enumerate}

\subsubsection{Check 2: Verification of Kernel Tracing}
We recompiled xv6 and tested tracing on \texttt{grep hello README} targeting \texttt{read} (bit 5, mask $2^5 = 32$):

\begin{lstlisting}[style=bash]
$ trace 32 grep hello README
3: syscall read -> 1023
3: syscall read -> 965
3: syscall read -> 453
3: syscall read -> 0
\end{lstlisting}
The output confirms that \texttt{README} was read in three chunks ($1023 + 965 + 453 = 2441$ bytes, matching the file's exact size), followed by a return code of \texttt{0} signaling EOF.

% ==============================================================================
\subsection{Part 3. Play with the mask}

\subsubsection{Bitmask Mechanics}
The trace mask is interpreted as a binary bitfield where bit position $N$ corresponds to syscall number $N$. Setting bit $N$ to 1 enables logging for that specific call.

\subsubsection{Experimental Sessions in QEMU}

\begin{enumerate}[leftmargin=*]
    \item \textbf{Tracing All Calls (Mask $2147483647 = 2^{31}-1$, all bits set):}
\begin{lstlisting}[style=bash]
$ trace 2147483647 grep hello README
4: syscall trace -> 0
4: syscall exec -> 3
4: syscall open -> 3
4: syscall read -> 1023
4: syscall read -> 965
4: syscall read -> 453
4: syscall read -> 0
4: syscall close -> 0
\end{lstlisting}
    Every phase of process setup, file access, and cleanup is intercepted.

    \item \textbf{Tracing Only Fork (Mask $2 = 2^1$, bit 1 for SYS\_fork):}
\begin{lstlisting}[style=bash]
$ trace 2 usertests forkforkfork
usertests starting
test forkforkfork: 3: syscall fork -> 4
4: syscall fork -> 5
5: syscall fork -> 6
5: syscall fork -> 7
5: syscall fork -> 8
6: syscall fork -> 9
5: syscall fork -> 10
6: syscall fork -> 11
5: syscall fork -> 12
6: syscall fork -> 13
5: syscall fork -> 14
6: syscall fork -> 15
5: syscall fork -> 16
6: syscall fork -> 17
6: syscall fork -> 18
5: syscall fork -> 19
6: syscall fork -> 20
5: syscall fork -> 21
5: syscall fork -> 22
6: syscall fork -> 23
5: syscall fork -> 24
6: syscall fork -> 25
5: syscall fork -> 26
6: syscall fork -> 27
5: syscall fork -> 28
5: syscall fork -> 29
6: syscall fork -> 30
6: syscall fork -> 31
5: syscall fork -> 32
5: syscall fork -> 33
5: syscall fork -> 34
5: syscall fork -> 35
5: syscall fork -> 36
6: syscall fork -> 37
5: syscall fork -> 38
6: syscall fork -> 39
6: syscall fork -> 40
7: syscall fork -> 41
7: syscall fork -> 42
5: syscall fork -> 43
5: syscall fork -> 44
5: syscall fork -> 45
7: syscall fork -> 46
5: syscall fork -> 47
5: syscall fork -> 48
6: syscall fork -> 49
5: syscall fork -> 50
5: syscall fork -> 51
6: syscall fork -> 52
6: syscall fork -> 53
5: syscall fork -> 54
6: syscall fork -> 55
5: syscall fork -> 56
5: syscall fork -> 57
6: syscall fork -> 58
6: syscall fork -> 59
5: syscall fork -> 60
5: syscall fork -> 61
5: syscall fork -> 62
5: syscall fork -> 63
6: syscall fork -> 64
5: syscall fork -> -1
OK
ALL TESTS PASSED
\end{lstlisting}
    Notice that multiple distinct PIDs (3, 4, 5, 6, 7) emit trace entries as child processes inherit the mask.

    \item \textbf{Tracing Only Write (Mask $65536 = 2^{16}$, bit 16 for SYS\_write):}
\begin{lstlisting}[style=bash]
$ trace 65536 echo hi
hi5: syscall write -> 2

5: syscall write -> 1
\end{lstlisting}
    The trace line is interleaved directly between the word \texttt{"hi"} and the terminating newline.
\end{enumerate}

\subsubsection*{Observe:}
\begin{enumerate}[leftmargin=*]
    \item \textbf{In trace 32 grep hello README why do you see only read, and not open or close?}\\[0.2em]
    The mask value \texttt{32} corresponds in binary to $2^5$ ($100000_2$), meaning only bit 5 is set. In xv6, \texttt{SYS\_read} is assigned index 5. The bitwise condition:
    \begin{center}
        \texttt{if ((p->trace\_mask >> num) \& 1)}
    \end{center}
    evaluates to true (\texttt{1}) if and only if \texttt{num == 5}. For \texttt{SYS\_open} (15) and \texttt{SYS\_close} (21), bits 15 and 21 are \texttt{0} in the mask \texttt{32}, causing the kernel to suppress logging for all other operations.

    \item \textbf{In forkforkfork, three different pids print lines. Which line of code from Part 2 makes that possible?}\\[0.2em]
    The assignment in \texttt{kernel/proc.c} inside \texttt{fork()}:
    \begin{lstlisting}[style=c]
np->trace_mask = p->trace_mask;
    \end{lstlisting}
    When parent process 3 forks child 4, child 4 inherits \texttt{trace\_mask = 2}. As child process 4 subsequently forks process 5, and process 5 forks process 6, each descendant process retains active tracing on its own \texttt{fork()} invocations.

    \item \textbf{In trace 65536 echo hi the trace line is printed between "hi" and the new line. Why? (How many times does echo call write?)}\\[0.2em]
    \texttt{echo} executes \textbf{two separate \texttt{write()} system calls}:
    \begin{itemize}[noitemsep]
        \item First, it writes the argument string: \texttt{write(1, "hi", 2)}. The characters \texttt{"hi"} appear immediately on the console display, and upon syscall completion, the kernel outputs:
        \begin{center}
            \texttt{5: syscall write -> 2}
        \end{center}
        \item Second, it writes the trailing newline: \texttt{write(1, "\textbackslash n", 1)}. The newline moves the cursor to the next line, and the kernel then outputs:
        \begin{center}
            \texttt{5: syscall write -> 1}
        \end{center}
    \end{itemize}
    Because the kernel trace is printed immediately following the completion of each individual syscall, the first trace output is emitted before \texttt{echo} has called \texttt{write()} for the newline.

    \item \textbf{Which mask traces both open (15) and close (21), and nothing else?}\\[0.2em]
    To select only bits 15 and 21, we compute the bitwise OR of their respective powers of two:
    \[
    \text{mask} = (1 \ll 15) \mid (1 \ll 21) = 2^{15} + 2^{21} = 32768 + 2097152 = \mathbf{2129920}
    \]
    In hexadecimal representation, this value is \textbf{\texttt{0x208000}}.
\end{enumerate}

% ==============================================================================
\section{Conclusion}

In Laboratory Work 10, I modified the \textbf{xv6-riscv} operating system kernel to implement a native, bitmask-configurable system call tracing facility.

This work made concrete the mechanisms of hardware privilege transitions. We followed the flow of control from a user-space invocation of \texttt{trace(mask)} through the generated assembly stub in \texttt{usys.S}, where the syscall index is loaded into register \texttt{a7} and an \texttt{ecall} instruction transitions the processor into supervisor mode. Inside the kernel, we traced how the trap vector routes control to \texttt{syscall.c}, which dispatches the request to \texttt{sys\_trace()} to record the bitmask in the process control block (\texttt{struct proc}).

By implementing output logging inside \texttt{syscall()}, we observed that system calls can be dynamically monitored, timed, and filtered without modifying target application binaries. Furthermore, propagating the trace mask across \texttt{fork()} in \texttt{kernel/proc.c} demonstrated how process attributes and debugging contexts are inherited throughout process trees.

Contrasting this implementation with the Linux \texttt{strace} analysis performed in Laboratory 9 illuminated the architectural difference between external process tracing via \texttt{ptrace} context switches and in-kernel tracing hooks, completing our practical investigation of operating system syscall mechanisms.

\end{document}
